\documentclass[10pt,a4paper]{amsart}
\usepackage{amsmath,amssymb,amsthm,graphicx}
\usepackage{verbatim}
\newenvironment{reference}{\comment}{\endcomment}
%\newenvironment{reference}{}{}
\newenvironment{slogan}{\comment}{\endcomment}
\newenvironment{history}{\comment}{\endcomment}

% For commutative diagrams we use Xy-pic
\usepackage[all]{xy}

% We use 2cell for 2-commutative diagrams.
\xyoption{2cell}
% Theorem environments.
%
\theoremstyle{plain}
\newtheorem{theorem}[subsection]{Theorem}
\newtheorem{proposition}[subsection]{Proposition}
\newtheorem{lemma}[subsection]{Lemma}

\theoremstyle{definition}
\newtheorem{definition}[subsection]{Definition}
\newtheorem{example}[subsection]{Example}
\newtheorem{exercise}[subsection]{Exercise}
\newtheorem{situation}[subsection]{Situation}

\theoremstyle{remark}
\newtheorem{remark}[subsection]{Remark}
\newtheorem{remarks}[subsection]{Remarks}

\numberwithin{equation}{subsection}

% Macros
%
\def\lim{\mathop{\mathrm{lim}}\nolimits}
\def\colim{\mathop{\mathrm{colim}}\nolimits}
\def\Spec{\mathop{\mathrm{Spec}}}
\def\Hom{\mathop{\mathrm{Hom}}\nolimits}
\def\Ext{\mathop{\mathrm{Ext}}\nolimits}
\def\SheafHom{\mathop{\mathcal{H}\!\mathit{om}}\nolimits}
\def\SheafExt{\mathop{\mathcal{E}\!\mathit{xt}}\nolimits}
\def\Sch{\mathit{Sch}}
\def\Mor{\mathop{\mathrm{Mor}}\nolimits}
\def\Ob{\mathop{\mathrm{Ob}}\nolimits}
\def\Sh{\mathop{\mathit{Sh}}\nolimits}
\def\NL{\mathop{N\!L}\nolimits}
\def\CH{\mathop{\mathrm{CH}}\nolimits}
\def\proetale{{pro\text{-}\acute{e}tale}}
\def\etale{{\acute{e}tale}}
\def\QCoh{\mathit{QCoh}}
\def\Ker{\mathop{\mathrm{Ker}}}
\def\Im{\mathop{\mathrm{Im}}}
\def\Coker{\mathop{\mathrm{Coker}}}
\def\Coim{\mathop{\mathrm{Coim}}}

% Boxtimes
%
\DeclareMathSymbol{\boxtimes}{\mathbin}{AMSa}{"02}

%
% Macros for moduli stacks/spaces
%
\def\QCohstack{\mathcal{QC}\!\mathit{oh}}
\def\Cohstack{\mathcal{C}\!\mathit{oh}}
\def\Spacesstack{\mathcal{S}\!\mathit{paces}}
\def\Quotfunctor{\mathrm{Quot}}
\def\Hilbfunctor{\mathrm{Hilb}}
\def\Curvesstack{\mathcal{C}\!\mathit{urves}}
\def\Polarizedstack{\mathcal{P}\!\mathit{olarized}}
\def\Complexesstack{\mathcal{C}\!\mathit{omplexes}}
% \Pic is the operator that assigns to X its picard group, usage \Pic(X)
% \Picardstack_{X/B} denotes the Picard stack of X over B
% \Picardfunctor_{X/B} denotes the Picard functor of X over B
\def\Pic{\mathop{\mathrm{Pic}}\nolimits}
\def\Picardstack{\mathcal{P}\!\mathit{ic}}
\def\Picardfunctor{\mathrm{Pic}}
\def\Deformationcategory{\mathcal{D}\!\mathit{ef}}
\usepackage{cleveref}
\crefname{originalreference}{original reference}{original references}
\Crefname{originalreference}{Original reference}{Original references}
\makeatletter
\newlabel{remark-top-chern-class}{{0FA4}{source}{}{sourcepointer-1}{}}
\newlabel{remark-top-chern-class@cref}{{[originalreference][0][]0FA4}{[1][0][]source}{}{}{}}
\makeatother
\crefname{originallemma}{original lemma}{original lemmas}
\Crefname{originallemma}{Original lemma}{Original lemmas}
\makeatletter
\newlabel{lemma-construction-gysin}{{0FBK}{source}{}{sourcepointer-2}{}}
\newlabel{lemma-construction-gysin@cref}{{[originallemma][0][]0FBK}{[1][0][]source}{}{}{}}
\makeatother
\crefname{originallemma}{original lemma}{original lemmas}
\Crefname{originallemma}{Original lemma}{Original lemmas}
\makeatletter
\newlabel{lemma-gysin-decompose}{{0FBL}{source}{}{sourcepointer-3}{}}
\newlabel{lemma-gysin-decompose@cref}{{[originallemma][0][]0FBL}{[1][0][]source}{}{}{}}
\makeatother
\crefname{originallemma}{original lemma}{original lemmas}
\Crefname{originallemma}{Original lemma}{Original lemmas}
\makeatletter
\newlabel{lemma-gysin-at-infty}{{0F9J}{source}{}{sourcepointer-4}{}}
\newlabel{lemma-gysin-at-infty@cref}{{[originallemma][0][]0F9J}{[1][0][]source}{}{}{}}
\makeatother
\crefname{originallemma}{original lemma}{original lemmas}
\Crefname{originallemma}{Original lemma}{Original lemmas}
\makeatletter
\newlabel{lemma-easy-gysin}{{02TB}{source}{}{sourcepointer-5}{}}
\newlabel{lemma-easy-gysin@cref}{{[originallemma][0][]02TB}{[1][0][]source}{}{}{}}
\makeatother
\crefname{originallemma}{original lemma}{original lemmas}
\Crefname{originallemma}{Original lemma}{Original lemmas}
\makeatletter
\newlabel{divisors-lemma-quasi-regular-immersion}{{063M}{source}{}{sourcepointer-6}{}}
\newlabel{divisors-lemma-quasi-regular-immersion@cref}{{[originallemma][0][]063M}{[1][0][]source}{}{}{}}
\makeatother
\title{The fundamental computation of the class of the virtual normal sheaf against the fundamental class of a regularly embedded subvariety}
\author{The Stacks Project Authors}
\date{stacks-project:chow.tex:local-snapshot-20261012; GNU FDL 1.2 or later}
\begin{document}
\maketitle
\noindent\textit{Editorial note.} All delivered mathematical text is the author's own, copied from the pinned Stacks Project chapter file chow.tex; no statement or proof step was rewritten, reordered or removed.\par\medskip
\noindent\textit{Editorial note.} Preamble provenance: the pinned chapter file opens with an input of the shared preamble and carries no preamble text of its own, so the project's single Stacks-wide preamble staged at sources/stacks-chow/preamble.tex is excerpted verbatim as two spans, lines 14-18 (Xy-pic setup) and lines 154-217 (the native theorem declarations and mathematical macros). Nothing was added, removed or redefined; the wrapper class and the amsmath/amssymb/amsthm/graphicx base packages are supplied by the assembler.\par\medskip
\noindent\textit{Editorial note.} The counted statement and proof contain no cross reference, no citation and no included file, so the delivered body is self contained.\par\medskip
\noindent\textit{Editorial note.} Canonical location of the counted proposition: Stacks tag 0FBM (sources/stacks-chow/pinned-original-source.tex line 12076).\par\medskip
\noindent\textit{Allowed standard prerequisites (editorial).}\begin{itemize}
\item The class of the virtual normal sheaf and its base change behaviour. Same chapter, Lemma labelled lemma-construction-gysin, cited at proof lines 12087-12089. Allows X to be replaced by Y in the computation.
\item Decomposition of the class of the virtual normal sheaf. Same chapter, Lemma labelled lemma-gysin-decompose, cited at proof lines 12097-12099. Reduces to the case where the normal and conormal sheaves are identified.
\item The class C at infinity and the Gysin class of the divisor W\_infinity. Same chapter, Lemma labelled lemma-gysin-at-infty, invoked at proof lines 12106-12112, and Lemma labelled lemma-easy-gysin, cited at proof lines 12116-12118. Computes C cap [X]\_n as [W\_infinity]\_n.
\item The symmetric algebra structure on the conormal sheaf of a regular immersion. Divisors, Lemma labelled divisors-lemma-quasi-regular-immersion, cited at proof lines 12120-12122. Identifies the normal cone with the vector bundle of the conormal sheaf, giving p\textasciicircum{}*[Z]\_\{n - r\} = [C\_Z X]\_n.
\end{itemize}\medskip
\ifdefined\hypertarget\hypertarget{sourcepointer-1}{}\fi
\noindent\textit{Source reference.} Original reference 0FA4: sources/stacks-chow/pinned-original-source.tex, source lines 7677--7677. Same-chapter pointer into the author environment remark at source line 7677; invoked at proof line 12089 of the counted proof. It is retained only as a navigation pointer; every load-bearing step is inside the delivered span.\par\medskip
\ifdefined\hypertarget\hypertarget{sourcepointer-2}{}\fi
\noindent\textit{Source reference.} Original lemma 0FBK: sources/stacks-chow/pinned-original-source.tex, source lines 11912--11912. Same-chapter pointer into the author environment lemma at source line 11912; invoked at proof line 12090 of the counted proof. It is retained only as a navigation pointer; every load-bearing step is inside the delivered span.\par\medskip
\ifdefined\hypertarget\hypertarget{sourcepointer-3}{}\fi
\noindent\textit{Source reference.} Original lemma 0FBL: sources/stacks-chow/pinned-original-source.tex, source lines 11994--11994. Same-chapter pointer into the author environment lemma at source line 11994; invoked at proof line 12095 of the counted proof. It is retained only as a navigation pointer; every load-bearing step is inside the delivered span.\par\medskip
\ifdefined\hypertarget\hypertarget{sourcepointer-4}{}\fi
\noindent\textit{Source reference.} Original lemma 0F9J: sources/stacks-chow/pinned-original-source.tex, source lines 10026--10026. Same-chapter pointer into the author environment lemma at source line 10026; invoked at proof line 12106 of the counted proof. It is retained only as a navigation pointer; every load-bearing step is inside the delivered span.\par\medskip
\ifdefined\hypertarget\hypertarget{sourcepointer-5}{}\fi
\noindent\textit{Source reference.} Original lemma 02TB: sources/stacks-chow/pinned-original-source.tex, source lines 5170--5170. Same-chapter pointer into the author environment lemma at source line 5170; invoked at proof line 12112 of the counted proof. It is retained only as a navigation pointer; every load-bearing step is inside the delivered span.\par\medskip
\ifdefined\hypertarget\hypertarget{sourcepointer-6}{}\fi
\noindent\textit{Source reference.} Original lemma 063M: sources/stacks-chow/tags.txt, source lines 7836--7836. Cross-chapter standard prerequisite located in the project tag list staged in the local snapshot; invoked at proof line 12120 of the counted proof. It is retained only as a navigation pointer; every load-bearing step is inside the delivered span.\par\medskip
\begin{lemma}
\label{lemma-gysin-fundamental}
In the situation described just above assume $\dim_\delta(Y) = n$
and that $\mathcal{C}_{Y \times_X Z/Z}$ has constant rank $r$.
Then
$$
c(Z \to X, \mathcal{N}) \cap [Y]_n =
c_{top}(\mathcal{E}) \cap [Z \times_X Y]_{n - r}
$$
in $\CH_*(Z \times_X Y)$.
\end{lemma}

\begin{proof}
The bivariant class $c_{top}(\mathcal{E}) \in A^*(Z \times_X Y)$ was
defined in Remark \ref{remark-top-chern-class}.
By Lemma \ref{lemma-construction-gysin} we may replace $X$ by $Y$.
Thus we may assume $Z \to X$ is a regular closed immersion
of codimension $r$, we have $\dim_\delta(X) = n$, and we have
to show that $c(Z \to X, \mathcal{N}) \cap [X]_n =
c_{top}(\mathcal{E}) \cap [Z]_{n - r}$ in $\CH_*(Z)$.
By Lemma \ref{lemma-gysin-decompose} we may even assume
$\mathcal{N}^\vee \to \mathcal{C}_{Z/X}$ is an isomorphism.
In other words, we have to show
$c(Z \to X, \mathcal{C}_{Z/X}^\vee) \cap [X]_n = [Z]_{n - r}$ in $\CH_*(Z)$.

\medskip\noindent
Let us trace through the steps in the definition of
$c(Z \to X, \mathcal{C}_{Z/X}^\vee) \cap [X]_n$. Let
$b : W \to \mathbf{P}^1_X$
be the blowing up of $\infty(Z)$. We first have to compute
$C \cap [X]_n$ where $C \in A^0(W_\infty \to X)$ is
the class of Lemma \ref{lemma-gysin-at-infty}.
To do this, note that $[W]_{n + 1}$
is a cycle on $W$ whose restriction to $\mathbf{A}^1_X$ is
equal to the flat pullback of $[X]_n$. Hence $C \cap [X]_n$
is equal to $i_\infty^*[W]_{n + 1}$. Since $W_\infty$ is an
effective Cartier divisor on $W$ we have
$i_\infty^*[W]_{n + 1} = [W_\infty]_n$, see Lemma \ref{lemma-easy-gysin}.
The restriction of this class to the open $C_ZX \subset W_\infty$
is of course just $[C_ZX]_n$. Because $Z \subset X$ is regularly
embedded we have
$$
\mathcal{C}_{Z/X, *} = \text{Sym}(\mathcal{C}_{Z/X})
$$
as graded $\mathcal{O}_Z$-algebras, see
Divisors, Lemma \ref{divisors-lemma-quasi-regular-immersion}.
Hence $p : N = C_ZX \to Z$ is the structure morphism of the
vector bundle associated to the finite locally free module
$\mathcal{C}_{Z/X}$ of rank $r$. Then it is clear that
$p^*[Z]_{n - r} = [C_ZX]_n$ and the proof is complete.
\end{proof}

\end{document}
